% ===== Block B03: 原文 p.25–30 | §1.3 带点曲线的形变 | 译手 tr-ch1b =====

\subsection{带点曲线的形变}

% [原文 p.25]
设 $C$ 是亏格 $g$ 的紧黎曼曲面。$C$ 的形变，按定义，是复空间之间的固有光滑全纯映射
$\pi:X\to Y$，连同指定的点 $y\in Y$，使得 $\pi^{-1}(y)$ 同构于黎曼曲面 $C$（见 1.1.2）。
更方便的是把复解析族 $f:X\to Y$ 的概念推广成：底空间 $Y$ 是复空间，而不必是复流形。
% [Tutor Note: 正文两处「见 1.1.2」指本书第 1 章 1.1.2 小节（复解析族的定义），不是公式编号；
%  形变的「固有」与「光滑」两个条件缺一不可：固有保证纤维紧，光滑保证 $\dd\pi$ 满射（纤维是光滑曲线）。原文 p.25]

例如，若 $Y=\operatorname{Specan}\C[\eps]/(\eps^2)$，则称 $C$ 的形变为 $C$ 的
无穷小形变（或一阶形变）。黎曼曲面 $C$ 的无穷小形变由上同调群
$H^1(C,\Theta_C)$ 参数化，其中 $\Theta_C$ 是 $C$ 上全纯向量场芽的层。（见 1.1.2。）
% [Tutor Note: $\operatorname{Specan}\C[\eps]/(\eps^2)$ 是「解析谱」：二重数环
%  $\C[\eps]/(\eps^2)$ 对应一个单点复空间，其上的族即把 $C$ 沿 $\eps$ 方向作一阶移动，
%  故称无穷小（一阶）形变。$H^1(C,\Theta_C)$ 参数化一阶形变是 Kodaira–Spencer 理论的基本结论。原文 p.25]

更一般地，我们可以定义带第 $n$ 无穷小邻域的 $N$-点黎曼曲面的数据
\[
\mathfrak{X}^{(n)}
=(C;Q_1,Q_2,\dots,Q_N;\eta_1^{(n)},\eta_2^{(n)},\dots,\eta_N^{(n)})
\]
\begin{sloppypar}的形变。带第 $n$ 无穷小邻域的亏格 $g$ 的 $N$-点黎曼曲面的无穷小形变由上同调群
$H^1(C,\Theta_C(-(n+1)\sum_{j=1}^N Q_j))$ 参数化。若 $C$ 是奇异半稳定曲线，
则 $C$ 的形变定义为复空间之间的固有平坦全纯映射 $\pi:X\to Y$，连同指定的点
$y\in Y$，使得 $\pi^{-1}(y)$ 同构于曲线 $C$。此时，曲线 $C$ 的无穷小形变由
上同调群 $\operatorname{Ext}^1_{\cO_C}(\Omega^1_C,\cO_C)$ 参数化；而带第 $n$
无穷小邻域的稳定 $N$-点曲线
\end{sloppypar}
\[
\mathfrak{X}^{(n)}
=(C;Q_1,Q_2,\dots,Q_N;\eta_1^{(n)},\eta_2^{(n)},\dots,\eta_N^{(n)})
\]
的无穷小形变，则由上同调群
\[
\operatorname{Ext}^1_{\cO_C}
\Bigl(\Omega^1_C,\cO_C\bigl(-(n+1)\sum_{j=1}^N Q_j\bigr)\Bigr)
\]
参数化。

% [原文 p.26]
（例如见 [Ar]、[DM, Section 1]、[SGA7, Exposé VI, 6]。）这里 $\Omega^1_C$ 是
曲线 $C$ 上 Kähler 微分的层。在我们的情形下，可以把正合列
\[
0\longrightarrow I_C/I_C^2\longrightarrow\Omega^1_Y\otimes\cO_C
\longrightarrow\Omega^1_C\longrightarrow 0
\]
当作层 $\Omega^1_C$ 的定义，其中 $I_C$ 是形变 $\pi:X\to Y$ 中 $C=f^{-1}(y_0)$
的定义理想层。置
$\Theta_C=\underline{\Hom}_{\cO_C}(\Omega^1_C,\cO_C)$。于是有正合列
% [Tutor Note: 下划线 $\underline{\Hom}$、$\underline{\operatorname{Ext}}^1$ 表示层同态与层 Ext
%  （取值于层），与整体截面群 $H^0$、$H^1$ 相区别；$\Theta_C=\underline{\Hom}_{\cO_C}(\Omega^1_C,\cO_C)$
%  是 $C$ 上全纯切层（全纯向量场层），与 p.25 出现的 $\Theta_C$ 是同一对象，只是那里未写下划线。
%  正合列 $0\to I_C/I_C^2\to\Omega^1_Y\otimes\cO_C\to\Omega^1_C\to0$ 是余切序列的拉回，
%  用它定义 $\Omega^1_C$ 的好处是对奇异纤维也成立。原文 p.26]
\begin{equation}
\begin{gathered}
0\longrightarrow H^1\bigl(C,\Theta_C\bigl(-(n+1)\sum_{j=1}^N Q_j\bigr)\bigr)
\longrightarrow
\operatorname{Ext}^1_{\cO_C}
\Bigl(\Omega^1_C,\cO_C\bigl(-(n+1)\sum_{j=1}^N Q_j\bigr)\Bigr)\\
\longrightarrow H^0\bigl(C,\underline{\operatorname{Ext}}^1_{\cO_C}
(\Omega^1_C,\cO_C)\bigr)\longrightarrow 0 .
\end{gathered}
\tag{1.53}
\end{equation}
% [Tutor Note: (1.53) 里三种对象并排：两端是整体截面群（$H^1$ 与 $H^0$），中间是层
%  $\operatorname{Ext}^1$ 的整体截面群。它来自局部–整体谱序列在曲线（复维数 1）上的退化；读法是：
%  中间项参数化全部无穷小形变，左端参数化「保持奇点」的形变（即每根分支各自平滑、只让结点处
%  的模参数变化），右端是奇点处的局部形变，故 $\simeq\C^q$。原文未提谱序列。原文 p.26]

若稳定曲线 $C$ 有 $q$ 个二重点 $P_1,P_2,\dots,P_q$，则
\[
\underline{\operatorname{Ext}}^1_{\cO_C}(\Omega^1_C,\cO_C)_Q=
\begin{cases}
\C, & \text{若 } Q=P_j,\ i=1,2,\dots,q,\\[2mm]
0, & \text{其他}.
\end{cases}
\]
% [Tutor Note: 上式的条件印作「$Q=P_j,\ i=1,2,\dots,q$」，指标 $i$ 与 $P_j$ 的下标 $j$
%  不一致，按上下文应读作「$Q=P_j$，$j=1,2,\dots,q$」。照原文排，不改。
%  另注意 $\underline{\operatorname{Ext}}^1$ 是支撑在奇点集上的层：在光滑点处为 $0$，
%  故 $H^0$ 只由 $q$ 个二重点贡献，得 $\C^q$。原文 p.26]
于是
\[
H^0\bigl(C,\underline{\operatorname{Ext}}^1_{\cO_C}(\Omega^1_C,\cO_C)\bigr)
\simeq\C^q .
\]
$H^1(C,\Theta_C(-(n+1)\sum_{j=1}^N Q_j))$ 的每个元素都对应数据
\[
\mathfrak{X}^{(n)}
=(C;Q_1,Q_2,\dots,Q_N;\eta_1^{(n)},\eta_2^{(n)},\dots,\eta_N^{(n)})
\]
的一个\emph{保持奇点}的无穷小形变。

\begin{definition}[1.23]
数据 $(\pi:Y\to B;s_1,s_2,\dots,s_N;\eta_1,\eta_2,\dots,\eta_N)$ 称为亏格 $g$ 的
\emph{带形式邻域的 $N$-点稳定曲线（全纯）族}，若下列条件满足。
\begin{enumerate}[label=(\arabic*)]
  \item $Y$ 与 $B$ 是连通复流形，$\pi:Y\to B$ 是固有平坦全纯映射，且
        $s_1,s_2,\dots,s_N$ 是 $\pi$ 的全纯截面。
  \item 对每一点 $b\in B$，数据
        $(Y_b:=\pi^{-1}(b);s_1(b),s_2(b),\dots,s_N(b))$ 是亏格 $g$ 的稳定
        $N$-点曲线。
  \item $\eta_j$ 是 $\cO_B$-代数同构
\[
\eta_j:\widehat{\cO}/s_j(B):=\varprojlim_{n\to\infty}\cO_Y/I_j^n
\simeq\cO_B[[\xi]],
\]
        其中 $I_j$ 是 $s_j(B)$ 在 $Y$ 中的定义理想。
\end{enumerate}
\end{definition}
% [Tutor Note: $\widehat{\cO}/s_j(B)$ 是「沿 $s_j(B)$ 的完备化」$\widehat{\cO}_{Y/s_j(B)}$
%  的原文写法（下标写在斜杠后面），其定义就是右边的逆极限 $\varprojlim_n\cO_Y/I_j^n$；
%  它同构于 $\cO_B[[\xi]]$ 正是「形式邻域」的含义：$\xi$ 是沿纤维方向的局部坐标，
%  $I_j$ 是 $s_j(B)$ 在 $Y$ 中的定义理想。这里 $B$ 上的系数环是 $\cO_B$，故是 $\cO_B$-代数同构。原文 p.26]

类似地，我们定义带第 $n$ 形式邻域的亏格 $g$ 的稳定 $N$-点曲线族
\[
(\pi:Y\to B;s_1,s_2,\dots,s_N;
\tilde{\eta}_1^{(n)},\tilde{\eta}_2^{(n)},\dots,\tilde{\eta}_N^{(n)}),
\]
只需把第三个条件换成

% [原文 p.27]
(3$'$) $\tilde{\eta}_j^{(n)}$ 是 $\cO_B$-代数同构
\[
\tilde{\eta}_j^{(n)}:\cO_Y/I_j^{n+1}\simeq\cO_B[[\xi]]/I_j^{(n+1)} .
\]

对全纯映射 $f:D\to B$，我们可以用 $f$ 定义底空间为 $D$ 的带第 $n$ 无穷小邻域
（resp. 形式邻域）的 $N$-点稳定曲线族的拉回 $f^*\mathfrak{X}^{(n)}$
（resp. $f^*\mathfrak{X}$）。

\begin{definition}[1.24]
亏格 $g$ 的带第 $n$ 形式邻域的稳定 $N$-点曲线族
\[
\mathfrak{X}=(\pi^{(n)}:C^{(n)}\to B^{(n)};
s_1^{(n)},\dots,s_N^{(n)};
\tilde{\eta}_1^{(n)},\dots,\tilde{\eta}_N^{(n)})
\]
称为在点 $b\in B^{(n)}$ 处\emph{通用}（resp.\ \emph{万有}），若对
$\pi^{(n)-1}(b)=(C;Q_1,\dots,Q_N;\eta_1,\dots,\eta_N)$ 的任意形变
\[
\mathfrak{Y}=(\pi:X\to Y;s_1,\dots,s_N;\hat{\eta}_1,\dots,\hat{\eta}_N)
\]
连同指定的点 $y\in Y$，都存在从 $Y$ 中 $y$ 的邻域到 $B^{(n)}$ 的全纯映射
（resp.\ 唯一的全纯映射）$f$，使得拉回 $f^*\mathfrak{X}$ 在 $Y$ 中 $y$ 的某个
邻域内同构于 $\mathfrak{Y}$，且 $\dd f$ 在基点处被唯一确定。若族在 $B^{(n)}$
的每一点处都通用（resp.\ 万有），则称族 $\mathfrak{X}$ 为通用（resp.\ 万有）族。
\end{definition}
% [Tutor Note: 通用（versal）与万有（universal）的唯一差别是映射 $f$ 的唯一性：曲线有非平凡
%  自同构时，模空间只能「粗」，族只能通用；自同构群平凡时才可能万有（见定理 1.28 末句）。
%  另注意记号：定义 1.19 用 $s_j$ 记形式邻域，本节改用 $\eta_j$（resp.\ $\eta_j^{(n)}$），
%  而本定义中 $\mathfrak{Y}$ 又写成 $\hat\eta_j$，同一对象三种记号，均照原文排。原文 p.27]

下文主要考虑带形式坐标的稳定带点曲线族，它作为带点曲线族是通用的。因此本节
只考虑稳定带点曲线族。稳定带点曲线族的小平–Spencer 映射由下面的定理给出。
证明见 [TUY, p.\ 499–503]。

\begin{theorem}[1.25]
对亏格 $g$ 的 $N$-点稳定曲线族 $(\pi:Y\to B;s_1,\dots,s_N)$ 及每一点
$b\in B$，存在 $\C$-线性映射
\begin{equation}
\rho_b:T_bB\longrightarrow\operatorname{Ext}^1_{\cO_{Y_b}}
\Bigl(\Omega^1_{Y_b},\cO_{Y_b}\bigl(-\sum_{j=1}^N s_j(b)\bigr)\Bigr),
\tag{1.54}
\end{equation}
其中 $Y_b=\pi^{-1}(b)$。
\end{theorem}

$\C$-线性映射 $\rho_b$ 称为该族 $(\pi:Y\to B;s_1,\dots,s_N)$ 在点 $b$ 处的小平–Spencer
映射。定理 1.25 的层版本如下：

\begin{corollary}[1.26]
若 $(\pi:Y\to B;s_1,\dots,s_N)$ 是亏格 $g$ 的 $N$-点光滑曲线族，则小平–Spencer
映射 $\rho_s$ 诱导 $\cO_B$-模同态
\[
\rho:\Theta_B\longrightarrow R^1\pi_*\underline{\Hom}
\Bigl(\Omega^1_{Y/B},\Theta_Y\bigl(-\sum_{j=1}^N s_j(B)\bigr)\Bigr).
\]
\end{corollary}

族的通用性判据如下：

\begin{proposition}[1.27]
族
\[
(\pi:\cC\to\cB;s_1,s_2,\dots,s_N)
\]
是亏格 $g$ 的稳定 $N$-点曲线族，它在点 $b\in\cB$ 处通用，当且仅当小平–Spencer
映射
\[
\rho_b:T_b\cB\longrightarrow\operatorname{Ext}^1_{\cO_{C_b}}
\Bigl(\Omega^1_{C_b},\cO_{C_b}\bigl(-\sum_{j=1}^N s_j(b)\bigr)\Bigr)
\]
在 $b$ 处是同构。
\end{proposition}

% [原文 p.28]
下面证明稳定 $N$-点曲线的通用族存在。

\begin{theorem}[1.28]
对亏格 $g$ 的每个 $N$-点稳定曲线 $(C;Q_1,\dots,Q_N)$，存在族
\[
\mathfrak{F}=(\pi:\cC\to\cB;s_1,s_2,\dots,s_N)
\]
与指定点 $b\in\cB$，使得 $\pi^{-1}(b)$ 同构于 $(C;Q_1,\dots,Q_N)$，且族
$\mathfrak{F}$ 在 $b$ 处通用。此外，$\cC$ 与 $\cB$ 是复流形，族 $\mathfrak{F}$
在 $\cB$ 中 $b$ 的某个小邻域的每一点处都通用。若该 $N$-点稳定曲线的自同构群
平凡，则族 $\mathfrak{F}$ 在 $b$ 处万有。
\end{theorem}

\noindent\textit{证明.}\quad 该定理是形变理论（[Ar]、[Sc]、[SGA7]）的推论。
由于我们需要通用族的显式描述，这里给出构造第一无穷小邻域的通用族的方法。

由假设，曲线 $C$ 只有通常二重点。因此由形变理论，存在通用族
$\pi^{(\phi)}:\cC^{(\phi)}\to\cB^{(\phi)}$ 与指定点 $x\in\cB^{(\phi)}$，
使得 $C_x=\pi^{(\phi)-1}(x)\simeq C$，且小平–Spencer 映射
\[
\rho_x:T_x\cB^{(\phi)}\longrightarrow
\operatorname{Ext}^1_{\cO_{C_x}}(\Omega^1_{C_x},\cO_{C_x})
\]
是同构。（由于 $C$ 的自同构群可能非平凡，族 $\pi^{(\phi)}$ 在 $x$ 处可能只是
通用的，而不是万有的。）置
\[
\cB=\cC^{(\phi)N}\setminus
\Bigl(\bigcup_{i<j}\Delta_{ij}\cup\{\cC^{(\phi)N}\text{ 的奇点}\}\Bigr),
\]
其中
\[
\Delta_{ij}=\{(x_1,\dots,x_N)\in\cC^{(\phi)N}\mid x_i=x_j\}
\]
是第 $(i,j)$ 个对角。存在自然的全纯映射 $p:\cB\to\cB^{(\phi)}$。再置
\[
\cC=\cC^{(\phi)}\times_{\cB^{(\phi)}}\cB,
\]
\begin{sloppypar}并令 $\pi:\cC\to\cB$ 为到第二个因子的投影。按定义，
$(Q_1,\dots,Q_N)\in p^{-1}(x)$。置 $x_0=(Q_1,\dots,Q_N)\in\cB$，则
$\pi^{-1}(x_0)=C_x\times x_0\simeq C$。此外，我们可以定义全纯截面
\end{sloppypar}
\[
s_j:\cB\longrightarrow\cC
\]
为
\[
s_j((P_1,\dots,P_N))
=(P_j,P_1,\dots,P_N)\in\cC^{(\phi)}\times_{\cB^{(\phi)}}\cB .
\]
于是 $s_j(x_0)=(Q_j,x_0)$。容易证明
$\mathfrak{F}=(\pi:\cC\to\cB;s_1,\dots,s_N)$ 在 $\cB$ 的每一点处通用。
\hfill$\square$

% [原文 p.29]
\figomit{1.6}{有序四点射影直线的模空间：$B_2=\C^2\setminus\{4\text{ 条直线}\}$。}
% [Tutor Note: 图注印「$B_2=\C^2\setminus\{4$ 条直线$\}$」，但图中画出 5 条：$u_1=0$、$u_1=1$、
%  $u_2=0$、$u_2=1$ 与对角线 $u_2=u_1$。按例 1.30 的条件，$B_2$ 应挖去 $u_j=0,1$（$j=1,2$）
%  四条直线，再加上 $u_1\ne u_2$ 这条对角线；图注的「4 条」疑为未计对角线。原文 p.29]

设 $\mathfrak{F}(\pi:=\cC\to\cB;s_1,s_2,\dots,s_N)$ 是亏格 $g$ 的稳定 $N$-点曲线
的通用族。置
% [Tutor Note: 原文此句印作 $\mathfrak{F}(\pi:=\cC\to\cB;s_1,\dots,s_N)$，括号内多一个等号，
%  应读作 $\mathfrak{F}=(\pi:\cC\to\cB;s_1,\dots,s_N)$（与定理 1.28 的写法一致）。照原文排，不改。原文 p.29]
\begin{equation}
\Sigma=\{P\in\cC\mid
\dd\pi_P:T_P\cC\to T_{\pi(P)}\cB \text{ 不是满射}\},
\tag{1.55}
\end{equation}
\begin{equation}
D=\pi(\Sigma).
\tag{1.56}
\end{equation}
% [Tutor Note: $\Sigma$ 是 $\dd\pi$ 的秩小于 $\dim\cB$ 的点，也就是纤维出现奇点（结点）的点；
%  对半稳定族，$\Sigma$ 恰是「结点轨迹」，$D$ 是它在底空间中的像（判别轨迹：纤维退化的参数值）。
%  引理 1.29 说 $\Sigma$ 余维 2、$D$ 是正规交叉除子，这就是 §1.2 条件 (3) 的几何来源。原文 p.29]
集合 $\Sigma$ 称为该族的\emph{临界轨迹}，$D$ 称为该族的\emph{判别轨迹}。
下面的引理是带通常二重点的奇异曲线形变理论的推论。（例如见 [Ar]、
[DM, Section 1] 或 [SGA7, Exposé VI, 6]。）

\begin{lemma}[1.29]
对亏格 $g$ 的稳定 $N$-点曲线的通用族
$(\pi:\cC\to\cB;s_1,s_2,\dots,s_N)$，设 $2g-2+N\ge1$。则
\begin{enumerate}[label=(\arabic*)]
  \item 
\[
\dim\cB=3g-3+N,\qquad \dim\cC=3g-2+N;
\]
  \item 临界轨迹 $\Sigma$ 是 $\cC$ 中余维 $2$ 的光滑子簇；
  \item 判别轨迹 $D$ 是 $\cB$ 中具正规交叉的除子。
\end{enumerate}
\end{lemma}
% [Tutor Note: 维数公式 $\dim\cB=3g-3+N$ 即「$g$ 维曲线的模数 $3g-3$ 加上 $N$ 个标点」，
%  $\dim\cC$ 再大 1（纤维方向）；与 $D$ 是除子相容（$\Sigma$ 余维 2 在 $\cC$ 中，投下后 $D$ 余维 1）。
%  条件 $2g-2+N\ge1$ 排除 $\PP^1$ 上少于 3 点、椭圆曲线上 0 点等自同构群非平凡（或稳定曲线不存在）
%  的退化情形。原文 p.29]

\begin{example}[1.30]
设 $(\PP^1;Q_1,\dots,Q_N)$ 是 $N$-点射影直线（亏格 $0$ 的光滑曲线）。
由 $\PP^1$ 的自同构，可以把前三个点 $Q_1,Q_2,Q_3$ 映到点 $0,1,\infty$。
我们可以把 $\PP^1$ 看作 $\C\cup\{\infty\}$。设 $u$ 是 $\C$ 的整体坐标。于是
$(\PP^1;Q_1,\dots,Q_N)$ 同构于
$(\PP^1;0,1,\infty,u_1,\dots,u_{N-3})$。置
\[
B_m=\{(u_1,\dots,u_m)\in\C\mid
u_j\ne0,1,\infty,\ u_i\ne u_j,\ i\ne j\}.
\]
例如，$B_1=\PP^1\setminus\{0,1,\infty\}$，而 $B_2=\C^2$ 挖去图 1.6 中的直线。
\end{example}
% [Tutor Note: 上式的 $B_m$ 印作 $(u_1,\dots,u_m)\in\C$，应为 $\C^m$：$B_2$ 是 $\C^2$ 中的
%  开集，有两个坐标 $u_1,u_2$。照原文排，已在报告 ERRATA 中标出。原文 p.29]

% [原文 p.30]
\figomit{1.7}{亏格 $0$ 的稳定四点曲线的通用族。}

设 $\pi:\cC=\PP^1\times B_m\to B_m$ 为到第二个因子的投影。对任意点
$(u_1,\dots,u_m)\in B_m$，置
\[
s_k((u_1,\dots,u_m))=k,\quad k=0,1,\infty;
\]
\[
\sigma_i((u_1,\dots,u_m))=u_i,\quad i=1,2,\dots,m .
\]
注意 $s_k$ 与 $\sigma_i$ 是 $\pi$ 的全纯截面。于是 $B_m$ 是有序 $(m+3)$-点
射影直线的模空间，而 $(\pi:\cC\to B_m;s_0,s_1,s_\infty;\sigma_1,\dots,\sigma_m)$
是有序 $(m+3)$-点射影直线的通用族。把通用族延拓成亏格 $0$ 的稳定
$(m+3)$-点曲线的族，最早由 Terada [T] 在另一背景下研究。亏格 $0$ 稳定
$N$-点曲线模空间的详细研究由 Gerritzen、Herrlich 与 van der Put [GHP] 完成。
$m=1$ 时延拓容易描述。$B_1$ 的紧化 $\overline{B}_1$ 是 $\PP^1$。设
$\tilde C$ 是 $\PP^1\times\overline{B}_1$ 在点 $(0,0)$、$(1,1)$、$(\infty,\infty)$
处的爆破，$\tilde\pi:\tilde C\to\overline{B}_1$ 为自然全纯映射。则 $s_k$ 与
$\sigma_1$ 分别有自然延拓 $\tilde s_k$ 与 $\tilde s_1$，它们是 $\tilde\pi$ 的
全纯截面。于是
$(\tilde\pi:\tilde C\to\overline{B}_1;\tilde s_0,\tilde s_1,\tilde s_\infty,\tilde s_1)$
是有序稳定四点亏格 $0$ 曲线的通用族 (1.30)。$\tilde\pi$ 在紧化点 $0,1,\infty$
上的纤维是带通常二重点的亏格 $0$ 稳定 $4$-点曲线。此例将在 §5.4 中讨论。
% [Tutor Note: 原文此处为 $\tilde s_k$ 与 $\tilde\sigma_i$ 两个延拓（即 $s_k$ 与 $\sigma_1$ 的延拓），
%  后文四元组的末项同样是 $\tilde\sigma_1$。$\tilde s_k$ 是固定标记 $0,1,\infty$ 的延拓，
%  $\tilde\sigma_1$ 是移动标记 $u_1$ 的延拓，二者不可混：若把 $\tilde\sigma_1$ 误抄成 $\tilde s_1$，
%  四元组末项就与第二项重复、且丢掉移动标记的延拓。另外「(1.30)」是对例 1.30 的交叉引用，
%  不是公式编号。原文 p.30]

由 Knudsen [Kn]，用有序亏格 $0$ 稳定 $5$-点曲线给出的 $B_2$ 的紧化
$\overline{B}_2$ 正是我们的 $\tilde C$。坐标的显式描述请参阅 [T] 与 [GHP]。
% [Tutor Note: $\dim B_2=3\cdot0-3+5=2$，所以 $B_2$ 的紧化必须是二维的；而
%  $\tilde C$ 是 $\PP^1\times\overline{B}_1$（曲面）在三个点处的爆破，也是二维，这就是
%  「$\overline{B}_2$ 正是 $\tilde C$」的含义。Knudsen 的一般构造（用 $N$ 点稳定曲线紧化
%  $\overline{M}_{0,N}$）本块不展开，见 [Kn]。原文 p.30]
